% القسم: المناهج
% الفصل: الاستقراء
% المبحث: الاستقراء على الأعداد الطبيعية

\documentclass[../../../include/open-logic-section]{subfiles}

\begin{document}

\olfileid{mth}{ind}{inN}

\olsection{الاستقراء على~$\Nat$}

الاستقراء في أيسر صوره طريق لإثبات حكم يعم الأعداد الطبيعية. ومبناه
أننا إذا ابتدأنا بـ$0$، ثم زدنا~$1$ مرة بعد مرة، بلغنا كل عدد طبيعي.
فلإثبات حكم لكل عدد يكفينا: (1)~إثباته للعدد $0$، و(2)~بيان أن ثبوته
لعدد~$n$ يوجب ثبوته للعدد الذي يليه~$n+1$. ولْنرمز إلى قولنا «للعدد
$n$ الخاصية $P$» بـ$P(n)$، وإلى قولنا «للعدد~$k$ الخاصية $P$»
بـ$P(k)$، وعلى هذا القياس؛ فإثبات $P(n)$ لكل $n \in \Nat$
بالاستقراء مؤلف من:
\begin{enumerate}
\item برهان $P(0)$، و
\item برهان أنه، لكل~$k$، إذا صدق $P(k)$ صدق $P(k+1)$.
\end{enumerate}
ولإظهار وجه ذلك، افرض حصول (1) و(2). فالأولى تعطينا صدق $P(0)$.
ومع الثانية نعلم، بوجه خاص، أن صدق $P(0)$ يستلزم صدق $P(0+1)$، أي
$P(1)$. وإنما استُخرج هذا من القول العام «لكل~$k$، إذا كان $P(k)$
فإن $P(k+1)$» بإحلال~$0$ محل~$k$؛ فينتج $P(1)$ بالقياس الشرطي. ثم
نرجع إلى (2)، ونحل هذه المرة $1$ محل~$k$، فنحصل على: إذا كان
$P(1)$ فإن~$P(2)$. وقد ثبت~$P(1)$، فينتج بالقياس الشرطي~$P(2)$،
وهكذا. فلكل عدد~$n$ نبلغ، بعد تكرار الانتقال $n$ مرة،~$P(n)$. ومن ثم
يثبت مجموع (1) و(2)~$P(n)$ لأي $n \in \Nat$.

ولنعتبر مثالًا. نريد أن نعرف كم مجموعًا مختلفًا يمكن رميه بـ$n$ من
أحجار النرد. ولنبدأ، وإن بدت البداية غريبة، بـ$0$ من الأحجار. فإذا لم
يكن عندك حجر أصلًا لم يكن ثمة إلا مجموع واحد يمكن أن «ترميه»: لا نقاط،
ومجموعها~$0$؛ فعدد النتائج المختلفة~$1$. ومع حجر واحد، أي $n=1$،
توجد ست قيم من $1$ إلى~$6$. ومع حجرين يمكن بلوغ كل مجموع من $2$ إلى
$12$، وعددها $11$. ومع ثلاثة أحجار يمكن بلوغ كل عدد من $3$ إلى
$18$، وعددها $16$. وتدل الأعداد $1$ و$6$ و$11$ و$16$ على نسق؛ فهل
الجواب $5n+1$؟ ثم إن $5n+1$ هو، لا محالة، الحد الأقصى الممكن؛ إذ لا
يوجد إلا $5n+1$ عددًا بين $n$، وهو أصغر مجموع ترميه بـ$n$ من الأحجار
إذا ظهر على جميعها $1$، وبين $6n$، وهو أكبر مجموع إذا ظهر على جميعها
$6$.

\begin{thm}
  يمكن بـ$n$ من أحجار النرد رمي جميع القيم الممكنة وعددها $5n+1$ بين
  $n$ و$6n$.
\end{thm}

\begin{proof}
لتكن $P(n)$ قولنا: «يمكن رمي أي عدد بين $n$ و~$6n$ باستعمال $n$ من
أحجار النرد». وإقامة الاستقراء تقتضي إثبات:
\begin{enumerate}
\item \emph{أساس الاستقراء} $P(0)$، ومعناه أن عدم الحجر لا يجيز إلا
  رمي المجموع $0$.
\item \emph{خطوة الاستقراء}: لكل $k$، إذا صدق $P(k)$ صدق~$P(k+1)$.
\end{enumerate}

وتصح (1)، لأن للرمية الخالية من الأحجار مجموعًا واحدًا هو $0$، وهو
القيمة الوحيدة الواقعة بين $0$ و~$6\cdot0$.

ولإثبات (2) نفرض مقدمة الشرطية، وهي~$P(k)$؛ ويسمى ذلك
\emph{فرض الاستقراء}. ثم نستعمله لإقامة~$P(k+1)$. وموضع الدقة أن
نصف القيم الممكنة عند رمي $k+1$ من الأحجار بدلالة القيم الممكنة عند
رمي $k$ منها، مع نتيجة الحجر الزائد رقم $k+1$. ولا سبيل إلى إعمال
فرض الاستقراء إلا بهذا الربط.

ومقتضى فرض الاستقراء إمكان تحصيل كل عدد بين $k$ و~$6k$ باستعمال
$k$ من الأحجار. فإذا رمينا~$1$ بالحجر رقم $(k+1)$ زاد المجموع $1$؛
فنبلغ كل قيمة من $k+1$ إلى $6k+1$ بأن نرمي الأحجار $k$، ثم يظهر~$1$
على الحجر رقم $(k+1)$. وتبقى القيم من $6k+2$ إلى $6k+6$؛ وهذه تحصل
إذا ظهر $6$ على كل واحد من الأحجار $k$، ثم ظهر عدد من $2$ إلى $6$
على الحجر رقم $(k+1)$. فثبت أنه باستعمال $k+1$ من الأحجار يمكن رمي
كل عدد من $k+1$ إلى $6(k+1)$؛ أي ثبت~$P(k+1)$ بالاعتماد على
الفرض~$P(k)$، وهو فرض الاستقراء.
\end{proof}

وكثيرًا ما يُستعمل الاستقراء لإثبات حكم في متوالية، من أعداد أو
مجموعات، تكون هي نفسها معرّفة «استقرائيًّا»؛ أي يُعرّف الشيء رقم
$(n+1)$ بواسطة الشيء رقم $n$. ومن ذلك أن نعرّف المجموع~$s_n$ للأعداد
الطبيعية إلى~$n$ هكذا:
\begin{align*}
  s_0 & = 0\\
  s_{n+1} & = s_n + (n+1)
\end{align*}
وينتج من التعريف:
\begin{align*}
  s_0 & = 0,\\
  s_1 & = s_0 + 1 && = 1,\\
  s_2 & = s_1 + 2 && = 1 + 2 = 3\\
  s_3 & = s_2 + 3 && = 1 + 2 + 3 = 6, \text{ وهكذا.}
\end{align*}
ومن هنا نستطيع أن نثبت بالاستقراء $s_n = n(n+1)/2$.

\begin{prop}
  $s_n = n(n+1)/2$.
\end{prop}

\begin{proof}
  يلزم إثبات (1)~$s_0 = 0\cdot(0 + 1)/2$، و(2) أن
  $s_k = k(k+1)/2$ يستلزم $s_{k+1} = (k+1)(k+2)/2$. أما (1) فظاهرة.
  وفي (2) نفرض فرض الاستقراء: $s_k = k(k+1)/2$، ثم نستعمله لبيان
  $s_{k+1} = (k+1)(k+2)/2$.

  ولنحسب $s_{k+1}$. فالتعريف يعطي
  $s_{k+1} = s_k + (k+1)$، وفرض الاستقراء يعطي
  $s_k = k(k+1)/2$. نعوّض بهذه القيمة في المعادلة الأولى، ولا يبقى
  إلا حساب يسير في الكسور:
  \begin{align*}
    s_{k+1} & = \frac{k(k+1)}{2} + (k+1)\\
    & = \frac{k(k+1)}{2} + \frac{2(k+1)}{2}\\
    & = \frac{k(k+1) + 2(k+1)}{2}\\
    & = \frac{(k+2)(k+1)}{2}.
  \end{align*}
\end{proof}

والحاصل أن الاستقراء هو الطريق الظاهر متى كان الحكم في متوالية $a_n$
معرّفة استقرائيًّا. وقد يفيد وإن لم تكن المتوالية كذلك، كما في قيم
رميات النرد، متى أمكن وصل حالة~$k+1$ بحالة~$k$ على وجه معلوم.

\end{document}
